部分子分布函数（PDF）描述夸克和胶子在质子内部按纵向动量的分布。它们具有
普适性：同样的PDF出现在许多不同的散射过程中，并且是通过对撞机数据的全局
拟合在唯象上确定的。除了最低的几个Mellin矩之外，PDF一直抗拒\emph{从头}
计算。格点QCD可用于计算质子的许多性质，但它本质上是一种欧氏空间方法，
而PDF是在闵氏空间中通过沿光锥分离的夸克产生与湮灭算符（此处我们关注的是
夸克与反夸克的PDF，而非胶子PDF。）的矩阵元来定义的。Ji提出了一种可能的
解决方案~\cite{Ji:2013dva}：利用等时算符
\begin{equation}
  \cO_\Gamma(x,\xi,n) \equiv
  \bar\psi(x+\xi n) \Gamma W(x+\xi n,x)\psi(x),
\end{equation}
的矩阵元来计算\emph{准部分子分布}（quasi-PDF），其中 $\psi$ 与
$\bar\psi$ 沿方向 $n$ 在空间上相隔距离 $\xi$，并由一条直Wilson线 $W$
连接。在质子动量分量 $p\cdot n$ 趋于无穷大的极限下，可以通过匹配
公式~\cite{Xiong:2013bka,Alexandrou:2015rja} 从准PDF得到PDF。

格点准PDF计算的挑战之一是 $\cO_\Gamma$ 的重正化；它是一个非定域算符，
且已知具有幂次发散~\cite{Ishikawa:2016znu,Chen:2016fxx}。该算符也出现在
伪PDF（pseudo-PDF）这一相关方法中~\cite{Radyushkin:2017cyf,Orginos:2017kos}，
类似的非定域算符还被用于研究横向动量依赖的部分子分布
（TMD PDF）~\cite{Musch:2011er,Engelhardt:2015xja,Yoon:2017qzo}。
准PDF的最初几项格点
研究~\cite{Lin:2014zya,Alexandrou:2015rja,Chen:2016utp,Alexandrou:2016jqi}
并未包含完整的重正化，但它们使用经涂抹（smearing）的规范链接来构建
Wilson线，微扰论已证明这可以减小幂次发散~\cite{Ishikawa:2016znu}。

文献~\cite{Constantinou:2017sej}在单圈格点微扰论中研究了 $\cO_\Gamma$
的重正化，发现手征对称性破缺允许 $\cO_\Gamma$ 与
$\cO_{\{\slashed{n},\Gamma\}}$ 混合。文献~\cite{Yoon:2017qzo}也给出了
这种混合的数值证据。非微扰重正化的研究由
文献~\cite{Alexandrou:2017huk,Chen:2017mzz}开创，它们采用
Rome-Southampton方法~\cite{Martinelli:1994ty}在RI-MOM方案中得到复的、
依赖于 $\xi$ 的重正化因子（存在混合时则为矩阵）$Z(\xi)$，随后再通过
微扰转换~\cite{Constantinou:2017sej}到唯象学所用的 $\MSbar$ 方案。
该方法的问题在于：必须确定的是整个函数而非寥寥数个参数，而且在大的
$\xi$ 处进行的转换可能落在微扰论适用的区域之外。由于中间方案固定的参数
多于最小必需数目，这意味着关联函数中的非微扰信息被丢弃，只能在向
$\MSbar$ 转换时以微扰方式找回。

在本工作中，我们研究使用一个仅在直线 $x+\xi n$ 上定义的辅助标量、颜色
三重态场 $\zeta(\xi)$ 来简化 $\cO_\Gamma$ 的重正化。在此方法中，我们把
QCD中含有 $\cO_\Gamma$ 的关联函数替换为扩展理论 QCD+$\zeta$ 中含有局域
色单态双线性 $\phi\equiv \bar\zeta\psi$ 的关联函数。这一方法很久以前
曾在连续理论中被使用过~\cite{Craigie:1980qs,Dorn:1986dt}。

在格点上，目前我们把 $n$ 限制为沿某一条坐标轴，即 $n=\pm\hat\mu$，并采用
作用量
\begin{equation}
\begin{gathered}
  S_\zeta = a\sum_\xi\frac{1}{1+am_0}\bar\zeta(x+\xi n)\left[\nabla_n+m_0\right]\zeta(x+\xi n),\\
\nabla_n\equiv\begin{cases}
n\cdot\nabla^* = \nabla^*_\mu & \text{若 }n=+\hat\mu \\
n\cdot\nabla = -\nabla_\mu & \text{若 }n=-\hat\mu,
\end{cases}
\end{gathered}
\end{equation}
其中 $\nabla$ 与 $\nabla^*$ 是格点上前向与后向协变导数，$a$ 为格距。
这在固定规范背景下给出裸传播子
\begin{equation}
  \left\langle \zeta(x+\xi n)\bar\zeta(x)\right\rangle_\zeta = \theta(\xi)e^{-m\xi}W(x+\xi n,x),
\end{equation}
其中 $m=a^{-1}\log(1+am_0)$，而 $W$ 是连接 $x$ 与 $x+\xi n$ 的格点规范
链接的简单乘积。可以用涂抹过的规范链接构造 $W$，只需以同一组规范链接定义
$\nabla_n$。质量项不会被任何对称性禁止，它对应于一个 $O(a^{-1})$ 抵消项。
利用这一传播子，当 $m=0$ 且 $\xi>0$ 时我们得到关系
\begin{equation}
  \cO_\Gamma(x,\xi,n) = \left\langle \bar\phi(x+\xi n)\Gamma \phi(x) \right\rangle_\zeta.
\end{equation}
对于 $\xi<0$，要得到 $\cO_\Gamma$，我们反转 $\zeta$ 传播的方向，并使用
$\cO_\Gamma(x,\xi,n)=\cO_\Gamma(x,-\xi,-n)$。

除确定抵消项 $m_0$ 之外，我们还必须重正化局域复合算符 $\phi$。格点夸克
作用量可能破坏手征对称性，此时 $\phi$ 可与 $\slashed{n}\phi$ 混合。于是
我们得到重正化模式
\begin{equation}
  \phi_R = Z_\phi\left(\phi+r_\text{mix}\slashed{n}\phi\right),\quad
  \bar\phi_R = Z_\phi\left(\bar\phi+r_\text{mix}\bar\phi\slashed{n}\right).
\end{equation}
我们也可以用投影算符构成算符
$\phi^\pm\equiv\frac{1}{2}(1\pm\slashed{n})\phi$，它以
$Z_\phi^\pm\equiv Z_\phi(1\pm r_\text{mix})$ 对角地重正化。对
$\xi\neq 0$，这一模式导致重正化的 $\cO_\Gamma$ 具有如下形式：
\begin{equation}
\begin{gathered}
  \cO_\Gamma^R(x,\xi,n) = Z_\phi^2e^{-m|\xi|}\cO_{\Gamma'}(x,\xi,n),\\
\Gamma' = \Gamma + r_\text{mix}\sgn(\xi)\{\slashed{n},\Gamma\} + r_\text{mix}^2\slashed{n}\Gamma\slashed{n}.
\end{gathered}
\end{equation}
因此，$\cO_\Gamma$ 可以通过确定三个参数来重正化：线性发散的 $m$、对数
发散的 $Z_\phi$，以及有限的 $r_\text{mix}$。由于 $r_\text{mix}$ 是
$O(g^2)$ 的，在单圈阶这与文献~\cite{Constantinou:2017sej}的模式相同。
在 $\xi=0$ 处，$\cO_\Gamma$ 是发散结构不同的局域夸克双线性算符，应有
一个单独的重正化因子~\cite{Constantinou_privcomm}，可以用通常方法计算；
我们采用文献~\cite{Alexandrou:2015sea}的结果。我们还注意到，由于局域
算符 $\phi$ 并非味道单态，即使 $\cO_\Gamma$ 是味道单态，夸克准PDF与胶子
准PDF之间也不存在混合。在后一情形下，混合将发生在向PDF匹配的过程中。

如果我们取 $n=\hat t$ 并赋予 $\zeta$ 自旋自由度（它们不与作用量耦合），
那么 $S_\zeta$ 就成为格点上静态夸克的作用
量~\cite{Eichten:1989zv,Sommer:2010ic}。特别地，$\phi$ 与静态-轻夸克
双线性算符相联系，我们发现重正化因子之间的关系：$Z_V^\text{stat}=Z_\phi^+$
与 $Z_A^\text{stat}=Z_\phi^-$。静态夸克理论还告诉我们如何消除 $O(a)$
格点人为效应~\cite{Kurth:2000ki}，即使格点上保持手征对称性，这类人为效应
依然存在。在连续理论中，$\cO_\Gamma$ 与静态夸克理论的关系此前已在
文献~\cite{Ji:2015jwa}中讨论过。

我们使用Rome-Southampton方法的一个变体，非微扰地确定这些重正化参数。
在 $N_f=4$ 扭转质量格点系综~\cite{Alexandrou:2015sea}上的朗道规范下，
我们计算位置空间的 $\zeta$ 传播子
\begin{equation}
  S_\zeta(\xi) \equiv \left\langle \zeta(x+\xi n)\bar\zeta(x)\right\rangle_{\text{QCD}+\zeta}
  = \bigl\langle W(x+\xi n,x)\bigr\rangle_\text{QCD},
\end{equation}
动量空间的夸克传播子 $S_\psi(p)$，以及 $\phi^\pm$ 的混合空间Green函数：
\begin{equation}
  G^\pm(\xi,p) \equiv
  \int d^4x\, e^{ip\cdot x}\left\langle
  \zeta(\xi n)\phi^{\pm}(0)\bar\psi(x)\right\rangle_{\text{QCD}+\zeta}.
\end{equation}
重正化参数以及 $\zeta$ 和 $\psi$ 场的重正化可由如下条件确定
\begin{gather}
  -\frac{d}{d\xi}\log\Tr S_\zeta(\xi)\Bigr|_{\xi=\xi_0} + m=0,\label{eq:cond_m}\\
\left[\frac{Z_\zeta}{3}\Tr S_\zeta(\xi_0)\right]^2=\frac{Z_\zeta}{3}\Tr S_\zeta(2\xi_0),\label{eq:cond_zeta}\\
\frac{1}{6}\frac{Z_\phi^\pm}{\sqrt{Z_\zeta Z_\psi}}\Re\Tr\left[S_\zeta^{-1}(\xi_0)G^\pm(\xi_0,p_0)S_\psi^{-1}(p_0)\right]=1,\label{eq:cond_phi}
\end{gather}
以及 $S_\psi$ 的RI$'$-MOM/RI-SMOM条
件~\cite{Martinelli:1994ty,Sturm:2009kb}。式~\eqref{eq:cond_m}对 $m$
敏感，而其余条件的构造则消除对 $m$ 的依赖。这些条件在标度 $\mu^2=p_0^2$
处定义了一族两参数的重正化方案，依赖于无量纲量 $y\equiv|p_0|\xi_0$ 与
$z\equiv p_0\cdot
n/|p_0|$。我们把这一族方案称为RI-xMOM。限制在运动学 $p_0\propto n$
（即 $z=1$）上，我们用维数正规化的微扰论~\cite{Green_forthcoming}计算了
到 $\MSbar$ 方案的单圈转换。在朗道规范下我们得到：
\begin{equation}
\begin{aligned}\label{eq:conversion}
   C_\phi &\equiv \frac{Z_\phi^{\MSbar}(\mu)}{Z_\phi^\text{RI-xMOM}(\mu,y,z=1)} \\
&= 1
+ \frac{\alpha_s C_F}{8\pi}\biggl[
  6\log\frac{y}{4}
  + 6\gamma_E - 8\log 2 + 7 \\
&\qquad\qquad - \cos y - \left(8\cos\frac{y}{2} - y\sin\frac{y}{2}\right)
\Ci\left(\frac{y}{2}\right) \\
&\qquad\qquad + 8\Ci(y)\biggr] + O(\alpha_s^2),
\end{aligned}
\end{equation}
其中 $\gamma_E$ 是 Euler-Mascheroni 常数，$\Ci$ 是余弦积分函数，
$\Ci(z) \equiv -\int_z^\infty
\frac{\cos(t)}{t} dt$。为了把 $m$ 转换到 $\MSbar$ 方案，我们采用
文献~\cite{Chetyrkin:2003vi,Melnikov:2000zc}中关于静态夸克传播子的三圈
结果。

\begin{figure}
  \centering
  \includegraphics[width=\columnwidth]{images/Eeff}
  \caption{裸辅助场传播子的有效能量，对应两种格距和三种不同的链接离散化。
    实心符号表示较细的格距，空心符号表示较粗的格距。曲线为三圈微扰结果，
    垂直平移 $-m$ 以便与较细格距上未涂抹的数据相匹配。其误差带表示
    $O(\alpha_s^3)$ 贡献的大小。}
  \label{fig:Eeff}
\end{figure}

在图~\ref{fig:Eeff}中，我们对两种不同格距：$a=0.082$~fm（$\beta=1.95$）
和 $a=0.064$~fm（$\beta=2.10$），展示式~\eqref{eq:cond_m}中的量\footnote{这是
  辅助场传播子的``有效能量''。文献~\cite{Musch:2010ka}曾研究过它，记作
  $Y_\text{line}$。}。加上 $m$ 即可使其重正化。不涂抹时，由于线性发散，
两个格距之间存在显著差异，但施加一步或五步HYP涂
抹~\cite{Hasenfratz:2001hp}后差异大幅减小，这同时也降低了大 $\xi$ 处的
统计不确定度。在小 $\xi/a$ 处，涂抹会扭曲形状，因此我们选择在形状相似的
$\xi_0\approx 0.6$~fm 处施加条件。随后我们在短距离处把细格距系综上使用
未涂抹链接的结果与微扰结果相匹配，从而转换到 $\MSbar$ 方案。

由式~\eqref{eq:cond_phi}可以分离出 $r_\text{mix}$ 的一个估计量。在我们的
数据中我们看到迹象表明它受到显著的 $O(a)$ 格点人为效应的影响。这些问题
本可通过Symanzik改进~\cite{Kurth:2000ki}来解决，但我们选择聚焦于不受
混合影响的螺旋度准PDF\footnote{按照推导自动 $O(a)$ 改进~\cite{Frezzotti:2003ni}
  所用的同一逻辑，也可以论证：对于处于最大扭转点的扭转质量格点QCD，混合
  算符对核子矩阵元的贡献将在 $O(a)$ 阶消失。}，它只依赖于
$r_\text{mix}^2$。我们发现后者并不比1\%大多少，在当前精度水平下可以
忽略。

\begin{figure}
  \centering
  \includegraphics[width=\columnwidth]{images/Zphi_ratio_vs_both}
  \caption{$\beta=2.10$ 的 $Z_\phi$，相对5HYP情形归一。左：固定
    $p\parallel n$ 与 $a^2p^2\approx 0.35$ 时随 $\xi$ 的变化。右：固定
    $\xi=6a$（未涂抹）与 $\xi=10a$（HYP）时随 $p^2$ 的变化。对每种链接
    离散化，我们取最小 $p^2$ 处两个值的平均。}
  \label{fig:Zratio_vs_both}
\end{figure}

剩余的因子 $Z_\phi$ 依赖于定义我们方案的运动学，但不同离散化的
$Z_\phi$ 之间的比值是与方案无关的。我们在大 $\xi$、小 $p$
（图~\ref{fig:Zratio_vs_both}）的平台区评估这个比值，以减小格点人为效应。
最后，我们使用未涂抹规范链接确定 $Z_\phi$，转换到 $\MSbar$ 方案并演化到
标度 2~GeV，如图~\ref{fig:Zphi}所示。我们发现单圈转换因子能有效地消除
对标度参量 $|p|\xi$ 的很大一部分依赖，两圈演化则消除了对标度 $|p|$ 的
大部分依赖。

\begin{figure}
  \centering
  \includegraphics[width=\columnwidth]{images/Zphi_scatter}
  \caption{$\beta=2.10$、使用未涂抹规范链接的 $Z_\phi$。数据显示了
    $p^2$ 与 $y\equiv |p|\xi$ 的一定范围；横轴为 $a^2p^2$，同一 $p^2$
    下不同 $y$ 略有错开。绿色空心方块给出我们的方案族中的值，蓝色实心
    菱形表示用式~\eqref{eq:conversion}在标度 $|p|$ 转换到 $\MSbar$
    的结果，带黑边的橙色实心三角形表示利用静态-轻夸克流的两圈反常量
    维~\cite{Chetyrkin:2003vi,Ji:1991pr,Broadhurst:1991fz}演化到标度
    2~GeV 后的 $\MSbar$ 结果。}
  \label{fig:Zphi}
\end{figure}

我们把上述重正化参数应用于在每个格距的一个 $N_f=2+1+1$ 扭转质量系综
计算的核子矩阵元。两个系综的物理参数相互匹配：$m_\pi\approx 370$~MeV
与 $p\cdot n\approx 1.85$~GeV。较粗的系综先前曾用于文
献~\cite{Alexandrou:2015rja,Alexandrou:2016jqi}，我们的方法论与
文献~\cite{Alexandrou:2016jqi}类似，包括在核子内插算符中使用动量涂
抹~\cite{Bali:2016lva}以便在大动量下获得良好信号。

\begin{figure*}
  \centering
  \includegraphics[width=0.495\columnwidth]{images/hel_comp}
  \includegraphics[width=0.495\columnwidth]{images/hel_comp_ren}
  \caption{$\beta=2.10$ 系综上的螺旋度准PDF矩阵元随 $\xi$ 的变化，三种
    不同的链接离散化，裸（左）与重正化（右）。仅显示 $\xi\geq 0$，因为
    实部关于 $\xi$ 为偶函数，虚部为奇函数。}
  \label{fig:ff_smearing}
\end{figure*}

\begin{figure}
  \centering
  \includegraphics[width=\columnwidth]{images/hel_comp_ens}
  \caption{两个系综上的螺旋度准PDF重正化矩阵元随 $\xi$ 的比较，均使用
    五步HYP涂抹。}
  \label{fig:ff_ens}
\end{figure}

图~\ref{fig:ff_smearing}展示了重正化对同位旋矢量螺旋度矩阵元
$\Delta h_{u-d}$ 的影响，
\begin{equation}
  \left\langle p,\lambda'\middle|\cO_{\slashed{n}\gamma_5}\middle|p,\lambda\right\rangle
\equiv\bar u(p,\lambda')\slashed{n}\gamma_5 u(p,\lambda) \Delta h(\xi,p\cdot n),
\end{equation}
它在细格距系综、不同链接离散化下计算。没有重正化时，不同链接类型之间
存在显著不一致，而重正化使它们很好地一致起来。在重正化矩阵元中，我们还
看到涂抹的好处：不涂抹时统计误差在大 $\xi$ 处迅速增大，$\xi\gtrsim 10a$
之后便没有可用信号。使用涂抹后，我们能够看到矩阵元在大 $\xi$ 处回归趋于
零。在大 $\xi$ 处，五步HYP涂抹也比一步给出更精确的数据。在
图~\ref{fig:ff_ens}中，我们比较了两个系综使用五步HYP涂抹的重正化数据。
二者符合得极好，这表明线性发散已被控制住，离散化效应不大。

\begin{figure}
  \centering
  \includegraphics[width=\columnwidth]{images/qpdf_hel_ren}
  \caption{$\beta=2.10$ 系综上的同位旋矢量螺旋度准PDF，三种不同链接离散
    化，由重正化矩阵元计算得到。}
  \label{fig:qpdf}
\end{figure}

最后，我们考察螺旋度准PDF，它由矩阵元的Fourier变换给出：
\begin{equation}
  \Delta\tilde q(x,p\cdot n) \equiv \frac{p\cdot n}{2\pi}\int d\xi\,
  e^{-ix\xi p\cdot n}\Delta h_q(\xi,p\cdot n).
\end{equation}
如文献~\cite{Alexandrou:2015rja}的图4所示，不重正化时未涂抹链接会导致
宽得多的分布。我们在图~\ref{fig:qpdf}中给出细格距系综上的重正化结果。
由于数据在大 $\xi$ 处噪声很大，我们把积分限制在 $|\xi|\leq 16a$
（未涂抹情形为 $|\xi|\leq 10a$）。当这种限制切掉部分信号时会引起振荡，
这在未涂抹情形中清晰可见。在未来的研究中，可以通过用一个描述矩阵元大
$\xi$ 行为的模型替代硬截断，或应用文献~\cite{Lin:2017ani,Chen:2017lnm}
提出的抑制大 $\xi$ 贡献的方法之一，获得改进的结果。忽略这些振荡，我们
看到重正化使不同链接离散化的数据达到了相当好的一致。

在这项工作中，我们证明了重正化格点准PDF这一非定域问题可以通过引入
辅助场转化为局域问题。辅助场方法也可以应用于横向动量依赖PDF格点研究中
使用的带订书钉形（staple-shaped）规范连接的算
符~\cite{Musch:2011er,Engelhardt:2015xja,Yoon:2017qzo}，那里的混合模式
将会不同。由于该方法与静态夸克理论紧密相连，可以利用该理论的已有结果，
例如文献~\cite{Chetyrkin:2003vi}的三圈连续计算。我们为下述结论提供了
支持证据：链接涂抹对这类计算是非常有益的技术，因为它使大 $\xi$ 处的
不确定性大幅降低；这一点早在静态夸克理论中就已被解
释~\cite{DellaMorte:2003mn}。从这里展示的结果到完整的PDF计算还有若干
步骤：从准PDF到PDF的匹配、对 $p\cdot n\to\infty$ 极限的控制、物理π介子
质量的使用，以及对有限体积效应与激发态效应的控制。

注：在我们中的一位在一次会议上报告本工作的初步版本之后~\cite{Green_Lat2017}，
我们得知存在一项基于辅助场方法的独立并行工作~\cite{Ji:2017oey}。我们也
注意到另一篇讨论准PDF可重正化性、但不使用辅助场方法的最新文
章~\cite{Ishikawa:2017faj}。
